\chapter{Desarrollo e implementación}
\label{chap:desarrollo-implementacion}

Documente la construcción, materialización y puesta en funcionamiento de los principales componentes desarrollados en el proyecto.

La estructura de este capítulo se adapta automáticamente a la modalidad seleccionada en el archivo \texttt{config.tex}.


% =============================================================
% MODALIDAD: INNOVACIÓN
% =============================================================

\ifdefstring{\ModalidadProyecto}{innovacion}{

\section{Solución desarrollada}
Describa la solución, producto, proceso o artefacto desarrollado, sus principales funcionalidades y la manera en que responde a la necesidad identificada.

\section{Arquitectura}
Presente la estructura de la solución, sus componentes principales y las relaciones entre ellos. Incluya, cuando corresponda, arquitectura técnica, modelo de datos, interfaces u otros elementos necesarios para comprender su funcionamiento.

\section{Desarrollo técnico}
Describa cómo se construyó y puso en funcionamiento la solución. Explique las principales decisiones técnicas, etapas, procedimientos y ajustes realizados durante su desarrollo.

\section{Integración}
Explique cómo la solución se articula con datos, usuarios, procesos, plataformas o sistemas existentes, cuando corresponda.

\section{Documentación técnica}
Presente los elementos necesarios para comprender, ejecutar, mantener o reutilizar la solución. Incluya, cuando corresponda, repositorios, instrucciones de instalación o ejecución, configuraciones, dependencias y documentación complementaria.

}{}


% =============================================================
% MODALIDAD: EMPRENDIMIENTO
% =============================================================

\ifdefstring{\ModalidadProyecto}{emprendimiento}{

\section{Propuesta de valor}
Explique el valor que el emprendimiento ofrece a sus clientes o usuarios, la necesidad que atiende y los principales elementos que diferencian la propuesta frente a otras alternativas.

\section{Producto o servicio}
Describa el producto o servicio desarrollado, sus características, atributos principales y forma de uso o prestación. Ajuste el título de esta sección cuando corresponda específicamente a un producto o a un servicio.

\section{Prototipo o MVP}
Presente el prototipo o producto mínimo viable desarrollado para poner a prueba la propuesta de valor. Describa su alcance, características y las principales decisiones tomadas durante su construcción. Ajuste el título de esta sección según el artefacto efectivamente desarrollado.

\section{Modelo de negocio}
Explique cómo el emprendimiento crea, entrega y captura valor. Incluya, según corresponda, segmentos de clientes, canales, relaciones con clientes, actividades, recursos, socios clave y fuentes de ingreso.

\section{Estructura financiera}
Presente los elementos financieros necesarios para comprender la lógica económica y la viabilidad del emprendimiento. Incluya, según corresponda, estructura de costos, estrategia de precios, fuentes de ingreso, proyecciones e indicadores financieros relevantes.

\section{Estrategia comercial}
Describa la estrategia utilizada para atraer, convertir y retener clientes. Incluya, cuando corresponda, acciones de tracción, marketing, comunicación, canales y ventas.

}{}